% Shared preamble of the two flow diagrams (flow_paper.tex, flow_x3.tex)
\usepackage{amsmath,amssymb,bm}
\usepackage[most]{tcolorbox}
\usetikzlibrary{positioning,arrows.meta,calc,fit,backgrounds}
\definecolor{cPlant}{HTML}{37474F}
\definecolor{cEnv}{HTML}{A0620F}
\definecolor{cCtrl}{HTML}{AD1457}
\definecolor{cAdapt}{HTML}{4527A0}
\definecolor{cX3}{HTML}{1565C0}
\definecolor{cOk}{HTML}{2E7D32}
\definecolor{cText}{HTML}{263238}
\definecolor{cRef}{HTML}{546E7A}
\newcommand{\vq}{\bm{q}}  \newcommand{\qv}{\bm{q}_v} \newcommand{\vw}{\bm{\omega}}
\newcommand{\vu}{\bm{u}}  \newcommand{\vd}{\bm{d}}   \newcommand{\vs}{\bm{s}}
\newcommand{\vtau}{\bm{\tau}} \newcommand{\veps}{\bm{\varepsilon}} \newcommand{\vF}{\bm{F}}
\newcommand{\vp}{\bm{p}}  \newcommand{\vv}{\bm{v}}
\newcommand{\that}{\hat{\theta}}
\newcommand{\sat}{\operatorname{sat}}
\newcommand{\R}{\mathbb{R}}
% \fblock{name}{position}{colour}{number}{title}{equations}{IN}{OUT}
\newcommand{\fblock}[8]{%
  \node[inner sep=0pt, outer sep=0pt] (#1) at #2 {%
  \begin{tcolorbox}[enhanced, width=8.4cm, colback=white, colframe=#3, colbacktitle=#3, coltitle=white,
     title={\sffamily\bfseries\small\tcbox[on line, colback=white, colframe=white, colupper=#3, arc=5pt,
     boxsep=0pt, left=2.5pt, right=2.5pt, top=0.5pt, bottom=0.5pt]{#4}\ #5}, boxrule=0.8pt, arc=3pt,
     left=5pt, right=5pt, top=4pt, bottom=4pt, colbacklower=#3!7, segmentation style={draw=#3!45},
     fontupper=\small, fontlower=\scriptsize\sffamily\color{cText}]
  \centering #6
  \tcblower\raggedright\textbf{IN}\enspace #7\par\smallskip\textbf{OUT}\enspace #8
  \end{tcolorbox}};}
\tikzset{
  sig/.style={-{Stealth[length=2.6mm]}, line width=0.9pt, color=cText},
  lbl/.style={font=\scriptsize\sffamily, text=cText, fill=white, inner sep=1.6pt, align=center},
  aux/.style={draw=cRef, fill=cRef!6, rounded corners=3pt, line width=0.7pt,
              font=\scriptsize\sffamily, text=cText, align=center, inner sep=5pt},
}
